\chapter{改訂履歴}
\label{app:revisions}

\section{版の履歴}
\label{sec:version-history}

本稿は建設中の理論であり、構成と内容は複数回にわたり作り替えられている。
主要な段階を記録する。各版の呼称は便宜的なものである。

版番号は段階に対応する。第 $n$ 段階が \texttt{v0.$n$.0} である。
構成の変更、命題の追加や撤回、新たな実証があれば段階を上げ、
誤記や体裁の修正は第三桁で扱う（\texttt{v0.8.1} など）。

\begin{center}
\begin{tabularx}{\textwidth}{@{}lllX@{}}
\toprule
段階 & 版 & 構成 & 主な内容 \\
\midrule
第1段階 & \texttt{v0.1.0} & 検討の時系列順 & $\Phi$ の定義、27類型、方法論の四分類 \\
第2段階 & \texttt{v0.2.0} & 論理順へ再編 & 理論・類型・方法・実証・総括の五部構成。証明と数値例を追加 \\
第3段階 & \texttt{v0.3.0} & 部構成の変更 & 制約下の演繹を「応用」から「写像の空間と制約」へ移動 \\
第4段階 & \texttt{v0.4.0} & 実証の追加 & 法人企業統計による $\CCC$ と $g^\star$ の測定 \\
第5段階 & \texttt{v0.5.0} & 弱点の類型化 & 独立性・外生性・水準と変化の三類型を整理 \\
第6段階 & \texttt{v0.6.0} & $\tau$ の導入 & $\CCC=\sum\tau_i+\tau_{\mathrm{inv}}$ により水準差を演繹 \\
第7段階 & \texttt{v0.7.0} & 既存理論の統合 & $\kappa_i$ の各項に既存理論を対応させ、付録\ref{app:priorwork}を三層化 \\
第8段階 & \texttt{v0.8.0} & 定義の補充 & 費用の分割、$E$、$\kappa_K$、$g^\star$ への $d$ と $I$ の導入 \\
\bottomrule
\end{tabularx}
\captionof{table}{版の履歴。各段階の版番号、構成の変更、主な内容。}
\label{tab:version-stages}
\end{center}

段階が進むごとに主張が弱くなる傾向がある。
第1段階で断定していた命題の多くは、
現在では仮定を明示した条件付きの形になっている。
これは理論の劣化ではなく、\textbf{適用範囲の特定}である。

\section{撤回した主張}

本稿は検討の過程で複数回にわたり主張を撤回している。
撤回した内容と理由を記録する。本文には訂正後の記述のみを残しているため、
どの主張が一度は成立していたと考えられていたかは、本付録にのみ現れる。

読者が途中まで読んだ段階で誤解しないよう、本文は訂正後の内容で一貫させてある。
以下は本稿がどのように誤ったかの記録であり、本文の理解には必要ない。

\section{測定に由来する訂正}

\begin{center}
\begin{tabularx}{\textwidth}{@{}XXl@{}}
\toprule
当初の記述 & 訂正後 & 契機 \\
\midrule
$\CCC$ の上昇により $g^\star$ が四半世紀で三分の一低下した
& $m$ が +91\% と分母の悪化 +49\% を上回り、$m/\CCC$ は $38.5\%\to49.2\%$ へ上昇
& $m$ の取得（命題\ref{prop:gstar-flat}） \\
\addlinespace
小規模層の制約は信用（$\kappa$）である
& 実際に拘束しているのは $m$。小規模層の $\CCC$ は大企業より短い
& 規模階層別の測定（命題\ref{prop:small-margin}） \\
\addlinespace
$\phi_{\rm cog}$ は世界のどこでも測定されていない
& 測定例が存在する。手法も確立している
& 先行研究の確認（第\ref{sec:phicog-measurable}節） \\
\addlinespace
小規模層の低い $m$ は固定費の分散不足による
& 役員報酬が主因。戻すと規模の序列が消える
& 人件費の内訳取得（命題\ref{prop:officer}） \\
\addlinespace
R2 では $\phi$ が目的関数から制約へ降格する
& 撤回。$H$ と $\beta$ の変化で説明でき、降格は内容を持たない
& 先行研究の確認（注\ref{rem:phi-demotion-retracted}） \\
\addlinespace
容量制約の時刻依存を理論に取り込む
& 撤回。$n=1$ の観測から27類型の基底を書き換えることになる
& 抽象と具体の非対称の確認（注\ref{rem:no-generalization}） \\
\addlinespace
予測Bの不良設定は筆者の予測設定の失敗である
& リスク・インセンティブ・トレードオフ論争として四半世紀未決着
& 先行研究の確認（第\ref{sec:predB-illposed}節） \\
\bottomrule
\end{tabularx}
\captionof{table}{測定に由来する訂正。当初の記述、訂正後の内容、契機。}
\label{tab:measurement-corrections}
\end{center}

このうち第一の訂正が最も重い。
$\CCC$ の悪化のみを観測して $g^\star$ の低下を結論しかけたが、
系\ref{cor:gstar}は比であり、分母の変化のみから結論は導けない。
$m$ を取得していれば避けられた誤りであり、
実際に本稿では「$m$ の変化を測定していないため解釈は保留する」と記して停止していた。
\textbf{保留したことが結果的に正しかった。}

\section{推論に由来する訂正}

\begin{center}
\begin{tabularx}{\textwidth}{@{}XX@{}}
\toprule
当初の記述 & 訂正後 \\
\midrule
在職中に信用が「継続的に売却されていた」
& 職務著作と同様の原始的帰属。個人に発生してから移転するのではない \\
\addlinespace
信用資産が最大となるのは退職直後である
& 最大となるのは信用ではなく変換原資（人脈・技術知識）。信用はゼロのまま \\
\addlinespace
$\lambda_{\rm novel}$ は各成分の和である
& 加法ではない。ロックファイルにより $\lambda_{\rm dep}$ は $\lambda_{\rm sec}$ に条件付く \\
\addlinespace
族5はデータ条件が良い
& 集計層では族2より弱い。個社層に降りて初めて成立する \\
\addlinespace
実データへのアクセス困難が検証の主要な障害である
& 族5の主要障害は（ii）仮説の不良設定。族2は（iv）アクセス制約だが、必要な粒度が公表されない特殊な形 \\
\bottomrule
\end{tabularx}
\captionof{table}{推論に由来する訂正。当初の記述と訂正後の内容。}
\label{tab:reasoning-corrections}
\end{center}

\section{事前登録が働いた記録}

以下は、事前に予測を文章として固定していなければ
事後的な解釈に流れていたと考えられる箇所である。

\begin{enumerate}[label=(\arabic*)]
\item \textbf{上限制手数料の消滅。}
族5において上限制が99.7\%消滅している事実は、
可測性仮説の傍証として解釈することが可能であった。
実際には上限制の価格水準が一桁低いだけで、可測性とは無関係である。

\item \textbf{$\CCC$ と売上高の同期負相関。}
45業種中31業種で負の相関（$p=0.008$）が得られた。
これを「$\CCC$ は景気指標として機能する」と読むことができたが、
分母効果と区別できていない。

\item \textbf{大企業の与信ポジション。}
$\mathrm{DSO}-\mathrm{DPO}$ が大企業で大きいことを
「大企業が取引先を支えている」と解釈しえた。
実際には予測Fが符号ごと棄却された結果である。
\end{enumerate}

第\ref{sec:narrative}節で他者の記述について述べた事後的物語化は、
筆者自身の解釈にも同様に生じる。

\section{仮定の明示化}

本稿の複数の命題は、当初は独立性の仮定を明示せずに述べていた。
検証の過程で、実際には同時に決まる量を独立として扱っていたことが判明し、
各命題に仮定を書き加えた。

\begin{center}
\begin{tabularx}{\textwidth}{@{}lXl@{}}
\toprule
命題 & 追加した仮定 & 契機 \\
\midrule
\ref{prop:pooling} & $\rho$ が $n$ に依存しない & 予測Jの検討 \\
\ref{prop:breakage} & $\dbar$ が外生に与えられる & 予測Kの検証（階段メニュー） \\
\ref{prop:capacity} & $m$ と $\E[\delta]$ が独立 & 同上 \\
\ref{prop:fcf} & $m$ と $\CCC$ が独立 & $m$ の取得（命題\ref{prop:gstar-flat}） \\
\bottomrule
\end{tabularx}
\captionof{table}{各命題に追加した独立性の仮定と、追加の契機。}
\label{tab:assumption-additions}
\end{center}

あわせて $\mathcal{G}$ の内生性（注\ref{rem:G-endogenous}）を追記した。
証明はいずれも変更していない。\textbf{適用範囲を狭めただけである}。

生産関数の定義域として未定義の記号 $\mathcal{X}$ を用いていた箇所を、
定義域を明示しない記述に改めた。
本稿は $f$ の内部構造を扱わないため、投入の空間を導入する必要がない。
$\mathcal{D}$、$\mathcal{F}$、$H$、$\phi_{\rm risk}$ を記号一覧に追加した。

\section{遅れ時間の導入}

第\ref{sec:tau}節の $\tau_i$ は、当初の版には存在しなかった。
第\ref{sec:level-change-gap}節で整理した弱点(C)、
すなわち水準と変化の未分離に取り組む過程で導入した。

導入の判断は、複数箇所で用いられるかを先に確認して行った。
水準差の説明（系\ref{cor:ccc-level}）、横断比較の不可能性（系\ref{cor:no-cross-comparison}）、
変化の分解（系\ref{cor:ccc-change}）の三つで用いる。

当初は成長期の $\CCC$ の歪みと $\beta$ 推定の再解釈も用途として見込んでいたが、
予測OとPの棄却により、年次データでは検出できないことが判明した
（第\ref{sec:lag-test}節）。この用途は導入の根拠から外した。

あわせて例\ref{ex:insolvency}が定義より前に $\CCC$ と $W$ を用いていた点を修正し、
$\kappa$ のみで記述する形に改めた。

\section{既存理論との関係の再評価}

本稿の草案では、$\kappa_i$ の各項について既存理論の有無を確認しないまま、
一部を本稿固有の領域と見込んでいた。文献を確認した結果、いずれにも理論が存在した。

\begin{center}
\begin{tabularx}{\textwidth}{@{}lXl@{}}
\toprule
当初の見込み & 確認された既存理論 & 分野 \\
\midrule
前受を取引信用として扱う視点は未開拓 & reverse trade credit \cite{mateut2014,daripa2011} & 企業金融 \\
労働の後払いに $\kappa$ の理論はない & bonding argument \cite{akerlof1986} & 労働経済学 \\
公開実績の蓄積は本稿固有 & 労働市場シグナリング \cite{spence1973} & 情報の経済学 \\
族2-4 の $\phi_{\rm cog}$ は未測定 & 前払い残高の退蔵の実証 \cite{nber2026pnbl} & 消費者行動 \\
\bottomrule
\end{tabularx}
\captionof{table}{$\kappa_i$ の各項について、当初の見込みと確認された既存理論、分野。}
\label{tab:prior-theory-check}
\end{center}

この確認により、寄与の所在が変わった。
当初は\emph{既存理論のない領域}を探していたが、
実際には\textbf{既存理論を各項の説明として導入し、
和として扱う部分が寄与}である（注\ref{rem:disjoint-literatures}）。
付録\ref{app:priorwork}を三層構造に改めたのはこのためである。

なお同種の失敗は繰り返されている。
第\ref{sec:predB-illposed}節で予測Bを立てた際にも、
リスク・インセンティブ・トレードオフ論争の存在を確認していなかった。

\section{定義の補充}

以下は誤りの訂正ではなく、定義が薄かった箇所の補充である。

\begin{center}
\begin{tabularx}{\textwidth}{@{}lXl@{}}
\toprule
対象 & 補充した内容 & 契機 \\
\midrule
生産関数 $f$ & 費用を $C_{\mathrm{prod}}+C_{\mathrm{admin}}$ に分割（定義\ref{def:cost}） & $\phi_{\mathrm{prod}}$ の源泉が未定義であった \\
$\phi_{\mathrm{prod}}$ & 限界費用と反実仮想価格で定義（定義\ref{def:phi-prod}） & 同上 \\
自己資本 $E$ & 払込資本と $\phi$ の累積として定義（式\eqref{eq:equity}） & 残余として暗黙であった \\
$\kappa_K$ & 残余請求権として定義（定義\ref{def:kappa-K}） & 出資を借入と同一に扱っていた \\
$g^\star$ & 減価償却 $d$ と投資 $I$ を含む形へ（系\ref{cor:gstar}） & $I\approx d\,r$ を暗黙に仮定していた \\
\bottomrule
\end{tabularx}
\captionof{table}{定義が薄かった箇所への補充内容と契機。}
\label{tab:definition-supplements}
\end{center}

$g^\star$ の改訂は測定値を変える。
1千万円未満の階層では、$m$ のみによる値 20.9\% に対し、
$I$ を含めた9年平均は 3.1\% となる。
命題\ref{prop:small-margin}の結論は維持されるが、
原因に投資負担が加わった。

費用の分割については、当初「$C$ は $\Phi$ に依存しない」を命題として証明していた。
しかし $C(\delta)$ と定義した時点で $\Phi$ が現れないのは定義の帰結であり、
証明は循環していた。撤回し、分割の定義として書き直したうえで、
$\Phi$ が費用を生む八つの経路を列挙した。

$\kappa_K$ の定義により、主観性が三段階に分かれることが判明した
（注\ref{rem:subjectivity}）。
評価写像 $v$ の主観性は情報の非対称であり、
残余請求権の価値 $V_j$ の主観性は確率測度の不一致である。
前者は開示で解消しうるが、後者は解消しない。

\section{可逆性の検討と訂正}

不可逆性の観点から枠組みを見直し、既存理論との照合を行った。
結果として一件の訂正と二件の追加が生じた。

\textbf{訂正。} 注\ref{rem:rho-endogenous}において、
$\rho$ が内生的になる機構を「相関の低い対象が枯渇する」と記述していた。
\cite{wagner2010} によれば正しい機構は異なる。
集約の対象が減るのではなく、\textbf{集約という行為そのものが集約者間の相関を生む}。
二主体が二リスクを等分に持てば、各主体は個別には安全になるが両者は完全に相関する。

\textbf{追加。} $\Phi$ の切り替え費用（注\ref{rem:switching-cost}）と
集約の社会的最適性（注\ref{rem:pooling-social}）を記録した。
前者は不可逆投資の理論、後者はシステミックリスクの研究に対応がある。

なお可逆性を費用の大小で分類する枠組み自体は、
不可逆投資の理論が既に持つものであり、本稿の寄与ではない。
決定が不可逆であるのは物理的に取り消せないためではなく
取り消しの費用が高いためである、という定式化は \cite{dixit1994} の系譜に属する。

\section{労働者視点の検討}
\label{sec:worker-predictions}

労働者を独立の対象として扱うかを検討し、扱わないと結論した。

当初は「労働者は一人事業の特殊例」という整理を試みた。
$\Phi$ が所与で顧客が一社に集中する一人事業、という見方である。
しかし $\Phi$ と $\phi$ の双方で制約が異なることが判明し、
この整理は撤回した（注\ref{rem:worker-constraints}）。
労働基準法が第\ref{sec:dof}節の六自由度のうち五つを制約しており、
$\phi_{\mathrm{prod}}$ と $\phi_{\mathrm{cog}}$ には対応物がない。

この検討の過程で四つの予測を登録した。

\begin{center}
\begin{tabularx}{\textwidth}{@{}lXl@{}}
\toprule
予測 & 内容 & 結果 \\
\midrule
Q & 資産の少ない労働者ほど支払サイクルが短い & 未検証 \\
R & 労働者は企業より高い $M_{\rm floor}$ を保持する & 未検証 \\
S & 雇用契約の類型の少なさは法の制約による & 未検証 \\
T & フリーランス新法が $\kappa$ の配分を移す & \textbf{検証不能} \\
U & 個人企業の営業利益率は法人の役員報酬を戻した値に近い & \textbf{支持} \\
\bottomrule
\end{tabularx}
\captionof{table}{労働者視点の検討で登録した予測と検証結果。}
\label{tab:worker-predictions}
\end{center}

予測Tの検証不能は、個人事業主の $\kappa$ が公表されないことによる。
この過程で第\ref{sec:obstacle-taxonomy}節の四分類に収まらない障害を発見した。
\textbf{調査はされるが集計されない}という形態である（注\ref{rem:tabulation-choice}）。

\section{障害の分類の拡張}

第\ref{sec:obstacle-taxonomy}節の四分類は、当初の調査経験に基づいて構成した。
その後の調査で、四分類に収まらない障害が二つ現れた。

\begin{center}
\begin{tabularx}{\textwidth}{@{}lXl@{}}
\toprule
種別 & 発見の契機 & 記録 \\
\midrule
（v）集計の選択 & 個人事業主の $\kappa$ を求めた際、調査票は存在するが集計表にない & 第\ref{sec:not-tabulated}節 \\
（vi）調整の不作動 & 予測Xの検定で、307社分のデータがあるにもかかわらず説明変数が動かない & 注\ref{rem:no-adjustment} \\
\bottomrule
\end{tabularx}
\captionof{table}{四分類に収まらなかった障害の種別、発見の契機、記録箇所。}
\label{tab:obstacle-extension}
\end{center}

（vi）は特に、本稿の方法論に対する含意が大きい。
第\ref{part:method}部は、検証できない場合の原因を
測定・理論・識別・アクセスの四つに帰していた。
\textbf{現実の側が理論の想定どおりに振る舞わない}という場合を想定していなかった。

なお、可測性仮説について予測Bを立てた段階では
$b$ が二値でしか取れないと記述していた。
その後の調査により手数料実績率という連続量が公表されていることが判明し、
検定が可能になった。障害はデータの粒度ではなかった。

\section{構成の変更}

内容の訂正とは別に、構成を二度変更した。

\begin{enumerate}[label=(\arabic*)]
\item 検討の時系列順から、理論・類型・方法・実証・総括の論理順へ。
\item 制約下の演繹（第\ref{ch:solo}章、第\ref{part:unmanned}章）を
独立した「応用」部から第\ref{part:catalog}部へ移動。
応用が正当化されるのは理論が実証を経た後であり、本稿の実証は完結していないため。
\end{enumerate}

第二の変更に伴い、64件の参照不整合（部と章の取り違え）と
2件の論理逆転（後続の章を過去形で参照）を修正した。
